\subsection{Limitations}
\label{sec:limitations}

We summarize what is established by the current evidence and then state the
main scope boundaries explicitly.

\paragraph{Established in this submission.}
The policy is runtime-legal and gold-free by construction
(Table~\ref{tab:feature_legality}); evaluation is disjoint by design across
the Aggregate-720 sources; and FTA leads each individual external baseline
(L1/S1/TALE) by point estimate with strictly positive source-stratified CI
lower bounds on Aggregate-720.
The pooled four-answer comparison is also reported and shows a closer gap
whose confidence interval includes zero on both evaluation sets.

\paragraph{Benchmark scope.}
Results are scoped to GSM8K (\texttt{openai/gsm8k}) under the matched-budget
protocol described in Section~\ref{sec:experiments}.
No accuracy claim is made beyond this dataset and task family.

\paragraph{Provider-specific scope.}
All main evidence is obtained on Cohere command-r-plus-08-2024.
The method is not claimed to generalize to other providers, model families,
or API versions without independent replication.

\paragraph{Seed~61 is failure-enriched.}
The 120-example seed~61 split was used as a development seed and is
failure-enriched: all four methods (frontier, L1, S1, TALE) fall to
approximately 54--58\% accuracy on this split, substantially below the
81--87\% range on seeds~41 and~71.
Absolute accuracy estimates from seed~61 alone are therefore not
representative of normal evaluation conditions.
Seed~61 is included in Aggregate-720 for aggregate coverage, and the
source-stratified confidence intervals account for its higher difficulty,
but it should not be interpreted as an independent IID validation split.

\paragraph{Pooled-ensemble comparison not statistically separated.}
Against a pooled four-answer majority-vote ensemble (Pooled-4), FTA leads by
point estimate on both evaluation sets ($+2.33$~pp on Final-300;
$+0.83$~pp on Aggregate-720), but the paired bootstrap confidence interval
includes zero in both cases: $[-0.67,\,+5.67]$ on Final-300 and
$[-1.11,\,+2.78]$ on Aggregate-720.
FTA cannot be claimed to be statistically superior to the pooled ensemble.
This comparison is reported as a robustness check, not as a retained
superiority claim.

\paragraph{Budget scope and compute accounting.}
The matched-budget claim is a per-method evaluation contract.
Full candidate-pool construction requires $4 \times B = 24$ logical inference
calls per example; FTA itself adds no post-generation model calls.
Monetary cost depends on provider pricing and is outside the main comparison
metric.

\paragraph{D9 multi-provider evidence is supporting only.}
The preliminary D9 experiment provides supporting evidence that the
failure-trace signal may generalize across providers, but the evaluation is
restricted to 80--160 examples per provider and uses cross-validation rather
than a held-out test set at full scale.
D9 findings are Tier~2 supporting evidence; they are not a second primary
submission claim.

\paragraph{Verifier-guided selection remains an open comparison.}
The manuscript cites verifier-guided and reward-model reranking work as a
related direction, but it does not compare FTA directly against trained
verifiers or lightweight learned rerankers. Such a head-to-head comparison is a
useful future direction, especially because FTA is intentionally deterministic,
auditable, gold-free, training-free, and judge-free at selection time.

\paragraph{Tie-breaker sensitivity.}
The TALE$\succ$S1$\succ$L1 fallback used when all three external answers
differ is deterministic and gold-free, but it is not claimed to be optimal.
Alternative tie orders move Aggregate-720 replay by less than one percentage
point in the available artifacts.

\paragraph{Residual pool-miss failures.}
A residual failure class exists in which all methods---frontier and external
baselines---produce incorrect answers for the same example.
The proposed policy cannot recover from this class, as no correct answer is
present in the candidate pool.

\paragraph{No formal risk-coverage objective.}
FTA does not optimize a formal risk-coverage objective and does not provide
formal risk-coverage guarantees; it is a deterministic rule-based policy that
uses runtime signals to decide when to defer, without guarantees on coverage or
selective risk \citep{geifman2017selectiveclassificationdeepneural,mozannar2023who}.

\paragraph{Not-retained variants.}
Follow-up prototypes were explored during development but did not meet
validation-grade evidence criteria.
They are reported as negative results in Table~\ref{tab:negative_results} and
are not part of the retained method.
